\section{Results}
\label{sec:results}

\subsection{Municipal aggregates of downscaled predictions}
Because pixel-level yield labels are unavailable, quantitative skill is assessed on municipal aggregates of the pixel predictions (Eq.~\ref{eq:mean-agg}).
Table~\ref{tab:main-metrics} summarizes validation and test metrics for the proposed spectral--climate downscaling model.
Validation $R^2$ peaks at epoch 12 ($0.227$), with RMSE $0.362$\,\si{t/ha} and MAPE $8.82\%$.
The corresponding test metrics are $R^2=0.228$, RMSE $0.351$\,\si{t/ha}, and MAPE $8.42\%$.
Absolute and relative municipal errors are therefore on the order of a few tenths of a tonne per hectare and about $8\%$--$9\%$ MAPE, with closely matched validation and test $R^2$ on the 2021 spatial holdout.

\begin{table}[t]
\centering
\caption{Validation and test metrics for municipal aggregates of the proposed downscaling model (best validation epoch 12). Units for RMSE/MAE are \si{t/ha}; MAPE is in percent.}
\label{tab:main-metrics}
\begin{tabular}{lcccc}
\toprule
Split & RMSE & MAE & $R^2$ & MAPE (\%) \\
\midrule
Validation (2021) & 0.362 & 0.287 & 0.227 & 8.82 \\
Test (2021)       & 0.351 & 0.271 & 0.228 & 8.42 \\
\bottomrule
\end{tabular}
\end{table}

Figure~\ref{fig:learning} shows learning dynamics.
Validation RMSE and $R^2$ improve through about epoch 12; after the scheduled learning-rate reduction near epoch 15, validation performance becomes unstable and eventually degrades, consistent with early stopping at epoch 22.

\begin{figure}[t]
\centering
\includegraphics[width=\textwidth]{learning_curves.pdf}
\caption{Training dynamics for the proposed spectral--climate downscaling model. The dashed line marks the best validation epoch (12).}
\label{fig:learning}
\end{figure}

\subsection{Architecture search without climate}
Table~\ref{tab:arch-top} lists the five best spectral-only configurations by validation $R^2$.
The top configuration is clearly ahead of the next-best models and is used as the spectral-only ablation below.
Wider spectral configurations often underperform, suggesting that capacity control and dropout matter for this dataset size and holdout design.

\begin{table}[t]
\centering
\caption{Top five spectral-only configurations by validation $R^2$ (24 completed runs). The proposed model instead uses the full Xavier climate layout with $d_{\mathrm{model}}=128$, four heads, $d_{\mathrm{inner}}=128$, and dropout $0.3$.}
\label{tab:arch-top}
\begin{tabular}{lcccccc}
\toprule
Rank & $d_{\mathrm{model}}$ & Heads & $d_{\mathrm{inner}}$ & Dropout & Val $R^2$ & Test $R^2$ \\
\midrule
1 & 96 & 8 & 64 & 0.3 & 0.165 & 0.117 \\
2 & 96 & 4 & 128 & 0.4 & 0.084 & 0.134 \\
3 & 64 & 4 & 64 & 0.4 & 0.076 & --- \\
4 & 64 & 4 & 128 & 0.3 & 0.071 & 0.119 \\
5 & 64 & 4 & 128 & 0.4 & 0.046 & 0.106 \\
\bottomrule
\end{tabular}
\end{table}

\subsection{Ablations}
\label{sec:ablations}

\paragraph{Spectral only.}
Training with spectral bands alone yields validation $R^2=0.165$ (RMSE $0.376$\,\si{t/ha}, MAPE $9.05\%$) and test $R^2=0.117$ (RMSE $0.375$\,\si{t/ha}, MAPE $9.09\%$).
Relative to the proposed spectral--climate model, removing climate covariates lowers both validation $R^2$ ($0.227\rightarrow 0.165$) and test $R^2$ ($0.228\rightarrow 0.117$), while MAPE remains near $9\%$.
Climate inputs therefore improve the municipal aggregates used for model selection and holdout evaluation under the 2021 spatial split.

\paragraph{Temporal pooling.}
With the best spectral-only architecture fixed, replacing NDVI-weighted temporal pooling by mean pooling over valid timesteps reduces validation $R^2$ from $0.165$ to $0.049$ (best epoch 12), while test $R^2$ remains nearly unchanged ($0.117$).
NDVI-weighted pooling therefore improves validation fit under the selection criterion used here, but does not translate into a clear test gain on this single holdout year.

\subsection{Per-year municipal diagnostics}
Figure~\ref{fig:scatter} shows predicted versus actual \emph{municipal} productivity for each harvest year under the proposed model.
Years 2020 and 2021 lie nearest the identity line (MAPE $8.1\%$--$8.6\%$), with 2021 $R^2=0.229$ consistent with the holdout evaluation.
Year 2022 remains difficult: RMSE rises to $1.14$\,\si{t/ha} and MAPE to $88.6\%$ despite near-zero $R^2$, reflecting a drought year with strong productivity collapse.
Years 2023--2024 are intermediate (Table~\ref{tab:per-year}).

\begin{figure}[t]
\centering
\includegraphics[width=\textwidth]{scatter_panel.pdf}
\caption{Predicted versus actual municipal soybean productivity (\si{t/ha}) by harvest year for the proposed downscaling model.
Points are colored by split (train / validation / test). The dashed line is identity.}
\label{fig:scatter}
\end{figure}

\begin{table}[t]
\centering
\caption{Per-year municipal metrics for aggregates of the proposed model (municipalities with coverage-filtered forecasts in each year).}
\label{tab:per-year}
\begin{tabular}{lccccc}
\toprule
Year & $n$ & RMSE & MAE & $R^2$ & MAPE (\%) \\
\midrule
2020 & 247 & 0.370 & 0.296 & 0.167 & 8.1 \\
2021 & 255 & 0.356 & 0.279 & 0.229 & 8.6 \\
2022 & 240 & 1.140 & 0.928 & $-0.001$ & 88.6 \\
2023 & 248 & 0.464 & 0.371 & 0.087 & 10.7 \\
2024 & 260 & 0.651 & 0.525 & 0.059 & 20.9 \\
\bottomrule
\end{tabular}
\end{table}

Figure~\ref{fig:map2021} maps municipal aggregate errors in the 2021 holdout year.
Figure~\ref{fig:guara} shows the corresponding intramunicipal product: pixel-level productivity contributions within Guarapuava, Paran\'{a}, obtained without field-level labels.
Spatial structure in such maps is the operational output of downscaling; municipal metrics only certify that the pixel field is consistent with PAM when averaged.

\begin{figure}[t]
\centering
\includegraphics[width=0.72\textwidth]{map_error_2021.png}
\caption{Municipal aggregate prediction error map for harvest year 2021 (holdout).}
\label{fig:map2021}
\end{figure}

\begin{figure}[t]
\centering
\includegraphics[width=0.72\textwidth]{guarapuava_heatmap.png}
\caption{Intramunicipal map of downscaled pixel-level productivity predictions for Guarapuava, Paran\'{a}.}
\label{fig:guara}
\end{figure}

\subsection{Early-season skill}
Table~\ref{tab:incomplete} and Figure~\ref{fig:incomplete} summarize early-season evaluation of municipal aggregates pooled over years.
With one or two season months, $R^2$ is negative and MAPE exceeds $34\%$.
Skill becomes clearly positive by month 3 ($R^2=0.266$) and peaks at month 5 ($R^2=0.429$, RMSE $0.648$\,\si{t/ha}, MAPE $26.2\%$), with a slight decline at month 6 ($R^2=0.411$).
Restricting the same protocol to 2021 alone, $R^2$ remains negative until month 6, where it reaches $0.229$ with RMSE $0.356$\,\si{t/ha}---matching the full-season holdout aggregate.
Early downscaled forecasts therefore improve as more of the season is observed when pooled across years, but the spatial holdout year requires nearly the full season window under this checkpoint.

\begin{table}[t]
\centering
\caption{Early-season evaluation of municipal aggregates (pooled years) as a function of season months $k$ retained from 1~October.}
\label{tab:incomplete}
\begin{tabular}{ccccc}
\toprule
$k$ & RMSE & MAE & $R^2$ & MAPE (\%) \\
\midrule
1 & 0.958 & 0.800 & $-0.247$ & 36.9 \\
2 & 0.907 & 0.793 & $-0.116$ & 34.6 \\
3 & 0.735 & 0.587 & 0.266 & 29.2 \\
4 & 0.680 & 0.523 & 0.373 & 27.3 \\
5 & 0.648 & 0.483 & 0.429 & 26.2 \\
6 & 0.658 & 0.477 & 0.411 & 26.9 \\
\bottomrule
\end{tabular}
\end{table}

\begin{figure}[t]
\centering
\includegraphics[width=\textwidth]{incomplete_series.pdf}
\caption{Early-season municipal metrics versus number of available season months $k$ (pooled years).}
\label{fig:incomplete}
\end{figure}
